% \documentclass[polish, 8pt]{article}
\documentclass[polish, 10pt]{extarticle}
\usepackage[top=2cm, bottom=2cm, left=3cm, right=3cm]{geometry}
\usepackage{graphicx} % Required for inserting images

\usepackage[polish]{babel}
\usepackage{polski}
\usepackage{placeins}
\usepackage{dirtytalk}
% Set page size and margins
% Replace `letterpaper' with `a4paper' for UK/EU standard size
\usepackage{csquotes}
% Useful packages
\usepackage{float}
\usepackage{amsmath}
\usepackage{graphicx}
\usepackage{hyperref}
\usepackage{graphicx}
\usepackage{listings}
\usepackage{xcolor}
\usepackage{amssymb}
\usepackage{enumitem}
\usepackage{float}
\definecolor{codegreen}{rgb}{0,0.6,0}
\definecolor{codegray}{rgb}{0.5,0.5,0.5}
\definecolor{codepurple}{rgb}{0.58,0,0.82}
\definecolor{backcolour}{rgb}{0.95,0.95,0.92}

\lstdefinestyle{mystyle}{
    backgroundcolor=\color{backcolour},   
    commentstyle=\color{codegreen},
    keywordstyle=\color{magenta},
    numberstyle=\tiny\color{codegray},
    stringstyle=\color{codepurple},
    basicstyle=\ttfamily\footnotesize,
    breakatwhitespace=false,         
    breaklines=true,                 
    captionpos=b,                    
    keepspaces=true,                 
    numbers=left,                    
    numbersep=5pt,                  
    showspaces=false,                
    showstringspaces=false,
    showtabs=false,                  
    tabsize=2
}
\lstset{style=mystyle}
\hypersetup{
    colorlinks=true,
    linktoc=all     %set to all if you want both sections and subsections linked
}
\usepackage{mdframed}
\usepackage{titling}
\setlength{\droptitle}{-2cm} 

\begin{document}

\title{Kolokwium I -- Zestaw A}
\author{}
\date{}

\maketitle
\vspace{-2cm}

\section*{Zadanie 1 (25 pkt)}
% W kapeluszu znajdują się dwie monety. Monety wyglądają tak samo, ale jedna jest uczciwa (z prawdopodobieństwem wypadnięcia orła równym \(1/2\)), natomiast druga jest stronnicza, z prawdopodobieństwem orła równym \(1/4\). Jedną z monet losowo wyciąga się z kapelusza, nie wiedząc, która to. Wybraną monetę nazywamy \emph{monetą \(C\)}.

% \begin{enumerate}
%   \item Monetę \(C\) rzucono dwukrotnie i za każdym razem wypadł orzeł. Znając tę informację, jakie jest prawdopodobieństwo, że moneta \(C\) jest uczciwa?
%   \item Czy zdarzenia „pierwszy rzut monety \(C\) daje orła” oraz „drugi rzut monety \(C\) daje orła” są niezależne? Uzasadnij.
%   \item Wyznacz prawdopodobieństwo, że w \(10\) rzutach monety \(C\) pojawią się dokładnie \(3\) orły. (Moneta jest jednakowo prawdopodobnie jedną z dwóch monet; nie zakładaj, że wcześniej wypadły dwa orły jak w punkcie (a).)
% \end{enumerate}
W kapeluszu znajdują się dwie monety, których nie można odróżnić wzrokiem. Jedna z nich jest uczciwa (prawdopodobieństwo orła $1/2$), a druga obciążona (prawdopodobieństwo orła $1/4$). Losujemy z kapelusza jedną monetę – nazwijmy ją monetą $C$.
\begin{enumerate}
\item Rzucono dwukrotnie monetą $C$ i za każdym razem wypadł orzeł. Jakie jest prawdopodobieństwo, że wylosowana moneta jest uczciwa?
\item Czy zdarzenia: „w pierwszym rzucie monetą $C$ wypadł orzeł” oraz „w drugim rzucie monetą $C$ wypadł orzeł” są niezależne? Uzasadnij odpowiedź.
\item Wracamy do momentu wyboru monety (zapominając o wynikach z punktu a). Jakie jest prawdopodobieństwo uzyskania dokładnie $3$ orłów w $5$ rzutach wylosowaną monetą $C$?
\end{enumerate}

\section*{Zadanie 2 (15 pkt)}
Niech \(A\) i \(B\) będą zdarzeniami losowymi.
Różnicą symetryczną \(A \triangle B\) nazywamy zdarzenie zachodzące wtedy, gdy zachodzi \(A\) lub zachodzi \(B\), ale nie oba jednocześnie (w logice odpowiada to alternatywie wykluczającej \emph{XOR}). Korzystając z definicji i własności prawdopodobieństwa, wykaż, że:
\begin{enumerate}
\item $P(A \triangle B) \;=\; P(A) + P(B) - 2P(A \cap B)$
\item $P(A) - P(B) \leq P(A \triangle B)$
\end{enumerate}

\section*{Zadanie 3 (10 pkt)}
Podaj definicję rozkładu prawdopodobieństwa zmiennej losowej. Sformułuj i udowodnij twierdzenie mówiące, że ten rozkład jest miarą probabilistyczną.

\section*{Zadanie 4 (30 pkt)}

Niech $X$ będzie zmienną losową o nieznanym parametrze $\theta \in \mathbb{R}_+$, której funkcja
gęstości prawdopodobieństwa ma postać
\[
f_\theta(x) =
\begin{cases}
\dfrac{2x}{\theta^2}, & \text{dla } 0 \leq x \leq \theta,\\[4pt]
0, & \text{w przeciwnym przypadku.}
\end{cases}
\]

\begin{enumerate}[label=\alph*)]
    \item Dla jakich wartości parametru $\theta$ funkcja $f_\theta$ jest
    funkcją gęstości prawdopodobieństwa zmiennej losowej $X$?

    \item Wyznacz dystrybuantę $F_X(x)$ zmiennej losowej $X$.

    \item Oblicz wartość oczekiwaną $E[2X+3]$, $Var(X)$ oraz medianę
    zmiennej losowej $X$ w zależności od parametru $\theta$.

    \item Oblicz prawdopodobieństwo $P\bigl(X^2 > 1\bigr)$ w zależności od parametru $\theta$.
\end{enumerate}

\section*{Zadanie 5 (20 pkt)}
Zmienna losowa \(X\) ma rozkład jednostajny \(U(-2, 3)\).
Wyznacz rozkład zmiennej losowej \(Y = (X-1)^2\).

\end{document}

